\documentclass[10pt,a4paper]{amsart}
\usepackage[margin=19mm]{geometry}
\usepackage{amsmath,amssymb,mathtools}
\usepackage[scr=rsfs,cal=euler]{mathalfa}
\usepackage{xfrac}
\usepackage[unicode,hidelinks,bookmarks=false]{hyperref}
\newcommand{\ie}{i.e.\ }
\newcommand{\cf}{cf.\ }
\newcommand{\resp}{respectively\ }
\newcommand{\Subsection}{\subsection}
\providecommand{\nobreakdash}{\nobreak\hbox{-}}
\setlength{\emergencystretch}{2em}
\multlinegap0pt
\usepackage{mathtools}
\usepackage{amssymb}
\usepackage{bm}
\usepackage{xcolor}
\newtheorem{lem}{Lemma}
\newtheorem{cor}{Corollary}
\newtheorem{prop}{Proposition}
\newcommand{\Z}{\mathbb{Z}}
\title{The rationality problem for multinorm one tori, II}
\author{Sumito Hasegawa, Kazuki Kanai, Yasuhiro Oki}
\date{Source: arXiv:2604.02848v1 (CC BY 4.0)}
\begin{document}
\maketitle

\noindent\emph{Source and licence.} The rationality problem for multinorm one tori, II. Version arXiv:2604.02848v1 (CC BY 4.0), distributed by arXiv under the Creative Commons Attribution 4.0 International licence, http://creativecommons.org/licenses/by/4.0/, as recorded in the arXiv licence metadata for this exact version and shown on the abstract page https://arxiv.org/abs/2604.02848v1. Full licence notice: \texttt{licenses/arxiv-2604.02848-CC-BY-4.0.txt}.\par\smallskip

\noindent\textit{Editorial scope.} Counted unit: the article's prop prop:cfcf (source0.tex lines 1459--1461). The article's complete original proof of it is delivered with it (source0.tex lines 1463--1496, one \texttt{proof} environment of 34 lines). No mathematics is written, repaired or re-derived: the statement and the proof are the article's own lines, in the article's own notation, and no retained line is substituted or re-flowed.  Retained uncounted as the source-local context the proof needs: lines 1273--1275; lines 1299--1304; lines 1327--1351; lines 1438--1449.  Nothing was omitted from any retained span, so the package carries no omission record.  No retained line contains a citation, so the wrapper carries no bibliography and no bibliography entry.  No retained line contains a figure environment, an image-inclusion command or drawing code, so no image is delivered and the package declares no asset dependency.  Editorial formatting only, in addition: an AMS \texttt{amsart} preamble that restates the article's own declarations and macros used by the retained lines, namely the environments \texttt{lem}, \texttt{cor}, \texttt{prop} on separate counters, together with the standard packages those lines need.  The retained environments print in this document as Lemma 1, Corollary 1, Proposition 1 in order of appearance, whereas the article prints the same items with the numbering of its own full document.  The complete native source of the article as distributed in the arXiv e-print for this exact version is bundled unchanged as \texttt{sources/197762/source0.tex}.
\begin{lem}\label{lem:zfsj}
The image of $P^{\langle \sigma_{2m}^{m}\rangle}$ under $\Psi_{m}$ contains $I_{m,0}$. 
\end{lem}

\begin{cor}\label{cor:eris}
We define $E_{m}$ as the kernel of $\Psi_{m}$. Then, there is an exact sequence of $D_{2m}$-lattices
\begin{equation}\label{eq:cdcf}
0\rightarrow E_{m} \rightarrow R_{m} \xrightarrow{\Psi_{m}} I_{m} \rightarrow 0. 
\end{equation}
\end{cor}

\begin{lem}\label{lem:rsdg}
We define an element $x_{1}$ of $R_{m}$ as follows: 
\begin{equation*}
x_{1}:=
\begin{cases}
\left(1+\sigma_{2m}^{m}\tau_{2m},(1+\tau_{2m})\sum_{j=1}^{(m-1)/4}(\sigma_{2m}^{(-1)^{j-1}(2j-1)}+\sigma_{2m}^{(-1)^{j-1}\cdot 2j}),-1\right)&\text{if }m\equiv 1\bmod 4;\\
\left(1+\sigma_{2m}^{m}\tau_{2m},-(1+\tau_{2m})\sum_{j=0}^{(m-3)/4}(\sigma_{2m}^{(-1)^{j}\cdot 2j}+\sigma_{2m}^{(-1)^{j}(2j+1)}),1\right)&\text{if }m\equiv 3\bmod 4. 
\end{cases}
\end{equation*}
\begin{enumerate}
\item We have $x_{1}\in E_{m}^{\langle \sigma_{2m}^{m}\tau_{2m} \rangle}$ and $\langle x_{1}\rangle_{\Z[D_{2m}]}\cong \Z[D_{2m}/\langle \sigma_{2m}^{m}\tau_{2m} \rangle]$. 
\item There is an isomorphism of $D_{m}$-lattices
\begin{equation*}
E_{m}\cong \Z[D^{(m)}/\langle \sigma_{2m}^{2}\rangle]\oplus \Z[D^{(m)}]
\end{equation*}
such that the composite
\begin{equation*}
\Z[D^{(m)}]\xrightarrow{c\mapsto (0,c)} \Z[D^{(m)}/\langle \sigma_{2m}^{2}\rangle]\oplus \Z[D^{(m)}]\cong E_{m}
\end{equation*}
is given by $1\mapsto x_{1}$. In particular, $E_{m}$ is permutation as a $D^{(m)}$-lattice, and 
\begin{equation*}
E_{m}=E_{m}^{\langle \sigma_{2m}^{2}\rangle}+\langle x_{1}\rangle_{\Z[D_{2m}]}. 
\end{equation*}
\end{enumerate}
\end{lem}

\begin{lem}\label{lem:flpp}
Let $G$ be a finite group, and
\begin{equation}\label{eq:abcf}
0\rightarrow E \rightarrow R \rightarrow M \rightarrow 0
\end{equation}
an exact sequence of $G$-lattices. Then it is a coflabby resolution of $M$ if and only if 
\begin{itemize}
    \item[(a)] $H^{1}(H,E)=0$; or
    \item[(b)] $R^{H}\rightarrow M^{H}$ is surjective,
\end{itemize}
for all subgroups of prime power orders in $G$. 
\end{lem}

\begin{prop}\label{prop:cfcf}
The sequence \eqref{eq:cdcf} is a coflabby resolution of $I_{m}$. 
\end{prop}

\begin{proof}
It suffices to prove (a) or (b) in Lemma \ref{lem:flpp} holds for any subgroup $H$ of prime power order in $D_{2m}$. Hence, we may assume that $H$ is contained in one of the following: 
\begin{equation*}
\langle \sigma_{2m}^{m}\rangle,\quad \langle \tau_{2m}\rangle,\quad \langle \sigma_{2m}^{m}\tau_{2m}\rangle,\quad \langle \sigma_{2m}^{m},\tau_{2m}\rangle,\quad \langle \sigma_{2m}^{2}\rangle. 
\end{equation*}

\textbf{Case 1.~$H=\langle \tau_{2m}\rangle$ or $H\subset \langle \sigma_{2m}^{2}\rangle$. }

In this case, $H$ is contained in $D^{(m)}:=\langle \sigma_{2m}^{2},\tau_{2m}\rangle$. Hence, Lemma \ref{lem:rsdg} implies that $E_{m}$ is permutation as an $H$-lattice. In particular, we obtain $H^{1}(H,E_{m})=0$ as desired. 

\textbf{Case 2.~$H=\langle \sigma_{2m}^{m}\rangle$. }

By direct computation, we have $I_{m}^{H}=I_{m,0}+\Psi_{m}(R_{m,2})$. Hence, Lemma \ref{lem:zfsj} gives an equality
\begin{equation*}
    I_{m}^{H}=\Psi_{m}(R_{m,1}^{H}\oplus R_{m,2}). 
\end{equation*}
Now, we obtain that (b) holds, that is, $R_{m}^{H}\rightarrow I_{m}^{H}$ is surjective. 

\textbf{Case 3.~$H=\langle \sigma_{2m}^{m},\tau_{2m}\rangle$. }

The equality $I_{m}^{H}=(1+\tau_{2m})(I_{m,0}+\Psi_{m}(R_{m,2}))$ can be confirmed directly. In particular, one has 
\begin{equation*}
    I_{m}^{H}=(1+\tau_{2m})I_{m}^{\langle \sigma_{2m}^{m}\rangle}. 
\end{equation*}
Hence, the surjection $\Psi_{m}\colon R_{m}^{\langle \sigma_{2m}^{m}\rangle}\rightarrow I^{\langle \sigma_{2m}^{m}\rangle}$, which is  ascertained in Case 2, implies the validity of (b). 

\textbf{Case 4.~$H=\langle \sigma_{2m}^{m}\tau_{2m}\rangle$. }

By the definition of $\Psi_{m}$, we obtain $I_{m}^{H}=(1+\sigma_{2m}^{m}\tau_{2m})(I_{m,0}+\Psi_{m}(R_{m,1}))$. In particular, we have
\begin{equation*}
    I_{m}^{H}=(1+\sigma_{2m}^{m}\tau_{2m})I_{m}. 
\end{equation*}
On the other hand, the homomorphism $R_{m}\rightarrow I_{m}$ is surjective, which is contained in Corollary \ref{cor:eris}. Therefore, we obtain that (b) is valid in this case.  
\end{proof}

\end{document}
